% Conclusiones (dictado del autor, con ajustes de forma).
% Estilo: no usar punto y coma (;) ni guion doble (--) en la redacción.

\section{Conclusiones}
\label{sec:conclusiones}

En este trabajo se evalúa la influencia del tipo de imagen y de la autosimilitud percibida sobre el rendimiento del FBC. Respecto a la pregunta de investigación, podemos determinar que para los casos evaluados sí se detecta una influencia de la autosimilitud en la calidad obtenida de la imagen. En cambio, la tasa y los tiempos de codificación y decodificación casi no dependen de la autosimilitud bajo las geometrías y el formato de codebook usados. La geometría de bloques, en cambio, tiene impacto sobre el tiempo de codificación, el de decodificación y la tasa de compresión obtenida. Frente a JPEG y PNG, podemos observar que en algunos casos la tasa de compresión del FBC fue más alta, pero a costa de perder calidad. Sin embargo, esos códecs aún presentan un mejor rendimiento en la relación entre tasa y distorsión y en el tiempo de codificación sobre los casos evaluados y frente a la implementación fractal realizada.

Como trabajo futuro se podría investigar otras implementaciones de la compresión fractal con bloques de tamaño variable dentro de cada imagen (hoy el tamaño es fijo por corrida), cambiando el tipo de almacenamiento del codebook y aplicando mejoras de cómputo con unidades gráficas. De este modo se podrían hacer pruebas más contundentes que pongan a prueba los códecs antes comparados.
